% !TEX root = ../main.tex

\begin{abstract}[zh]
\addcontentsline{toc}{chapter}{摘 \quad 要}
科学阅读并非一系列彼此孤立的检索：论文会随版本修订而变化，生成式解释依赖所选证据，研究者的兴趣也会在多次使用过程中持续演化。
现有的检索增强生成、引文评测与推荐研究分别覆盖了这一流程的部分环节，但尚未共同定义一种移动学术助手，使论文来源、生成主张和自适应用户状态都可被检查。
本文将 MNEME 定义为一个证据优先的移动文献闭环，连接版本感知的论文表示、受证据约束的主张级辅助、可由用户纠正的兴趣记忆，以及边界明确的移动交互层。
现有项目材料能够支撑可追溯的系统设计、交互式移动原型和已汇报的形成性观察；然而，产品源码的固定版本、E1--E3 实验输出以及 E4 的参与者级记录尚未提供，因此实现状态与有效性结果仍明确标记为待验证。
由此，本文当前的贡献是一套论证完整且可直接补充证据的系统规范：它规定了判断 MNEME 能否在不混淆流畅输出与已验证支持的前提下实现持续文献理解所需的不变量、评测方案与主张边界。
\end{abstract}

{
\newfontface{\arial}{Arial}[Scale=0.94]
\ctexset{chapter/format+={\arial}}
\begin{abstract}[en]
\addcontentsline{toc}{chapter}{ABSTRACT}
Scientific reading is not a sequence of isolated searches: papers change across revisions, generated explanations depend on selected evidence, and a researcher's interests evolve across sessions.
Current retrieval-augmented generation, citation evaluation, and recommendation research addresses parts of this workflow, but does not by itself define a mobile scholarly assistant whose paper provenance, generated claims, and adaptive user state remain jointly inspectable.
This thesis defines MNEME as an evidence-first mobile literature loop that connects revision-aware paper representation, claim-level evidence-bounded assistance, a user-correctable interest memory, and a bounded mobile interaction layer.
The supplied project materials establish a traceable system design, an interactive mobile prototype, and reported formative observations; however, the product source revision, E1--E3 experimental outputs, and participant-level E4 records were not supplied, so implementation and effectiveness results remain explicitly pending.
The resulting contribution is therefore an argument-complete and evidence-ready system specification: it states the invariants, evaluation protocols, and claim boundaries required to determine whether MNEME provides continuous literature understanding without confusing fluent output with verified support.
\end{abstract}
}
